%Royal Statistical Society Series B Statistical Methodology
\documentclass[12pt]{article}

%    \renewcommand{\baselinestretch}{1.1}

%Packages
\usepackage{authblk}
\usepackage{amsmath,amsfonts,amssymb,lmodern,geometry,enumerate}
%\usepackage{cancel}
\usepackage{mathtools} 
\usepackage[normalem]{ulem}
%\usepackage[justification=centering]{caption}
\usepackage[font=small,labelfont=bf]{caption}
%\usepackage{tkz-euclide}
%\usetkzobj{all}
%\usepackage{refcheck}
%\usepackage{showkeys}
%\usepackage{subcaption}
\usepackage[T1]{fontenc}
\usepackage[latin1]{inputenc}
\usepackage[english]{babel}
\usepackage{lmodern}
%\usepackage{scalefnt}
\usepackage{dsfont}
\usepackage{stmaryrd}
\usepackage{color}
\usepackage{bm,bbm}
\usepackage{mathrsfs,url}
\usepackage[hidelinks]{hyperref}
%\usepackage{breakcites} %for breaking the citations
%\usepackage{textcase}
\usepackage{graphicx}
\usepackage{wrapfig}
%\usepackage{wasysym}
%\usepackage{flushend,cuted}
\usepackage{bm}
\usepackage{tabularx}
\usepackage{indentfirst}
\usepackage{xparse}
\usepackage{tikz}
\usepackage{pgfplots}
\pgfplotsset{compat=newest}
%\usepackage{mdwlist}
%\usepackage{tkz-graph}
%\usepackage{transparent}
\usepackage{amsthm}
\usepackage{nameref}
%\GraphInit[vstyle = Shade]

\topmargin -1.5cm \evensidemargin 0.5cm \oddsidemargin 0.5cm
\textwidth15.8cm \textheight22cm
\parskip10pt


%Shortcuts for Theorems


\newtheorem{thm}{Theorem}[section]
\newtheorem{prop}[thm]{Proposition}
\newtheorem{defi}[thm]{Definition}
\newtheorem{lma}[thm]{Lemma}
\newtheorem{cor}[thm]{Corollary}
\newtheorem{exam}[thm]{Example}
\newtheorem{countexam}[thm]{Counterexample}
\newtheorem{rem}[thm]{Remark}
\newtheorem{con}[thm]{Conjecture}
\newtheorem{assum}[thm]{Assumption}

\newenvironment{Proof}{\removelastskip\par\medskip
\noindent{\em Proof.} \rm}{\penalty-20\null\hfill$\square$\par\medbreak}
\allowdisplaybreaks

\makeatletter
%%% Define shorter display equations commands
\def\be#1{\begin{equation*}#1\end{equation*}}
\def\ben#1{\begin{equation}#1\end{equation}}
\def\bes#1{\begin{equation*}\begin{split}#1\end{split}\end{equation*}}
\def\besn#1{\begin{equation}\begin{split}#1\end{split}\end{equation}}
\def\bea#1{\begin{align*}#1\end{align*}}
\def\bean#1{\begin{align}#1\end{align}}
\def\bg#1{\begin{gather*}#1\end{gather*}}
\def\bgn#1{\begin{gather}#1\end{gather}}
\setlength{\multlinegap}{1em}
\def\bml#1{\begin{multline*}#1\end{multline*}}
\def\bmln#1{\begin{multline}#1\end{multline}}

\def\note#1{\par\smallskip%
\noindent%
\llap{$\boldsymbol\Longrightarrow$}%
\fbox{\vtop{\hsize=0.98\hsize\parindent=0cm\small\rm #1}}%
\rlap{$\boldsymbol\Longleftarrow$}%
\par\smallskip}
\newcommand{\leftsup}[2]{{}^{#1}\! #2}%left upper index

%%% Define bracket commands
\def\given{\mskip 0.5mu plus 0.25mu\vert\mskip 0.5mu plus 0.15mu}
\newcounter{@bracketlevel}
\def\@bracketfactory#1#2#3#4#5#6{
\expandafter\def\csname#1\endcsname##1{%
\addtocounter{@bracketlevel}{1}%
\global\expandafter\let\csname @middummy\alph{@bracketlevel}\endcsname\given%
\global\def\given{\mskip#5\csname#4\endcsname\vert\mskip#6}\csname#4l\endcsname#2##1\csname#4r\endcsname#3%
\global\expandafter\let\expandafter\given\csname @middummy\alph{@bracketlevel}\endcsname
\addtocounter{@bracketlevel}{-1}}%
}
\def\bracketfactory#1#2#3{%
\@bracketfactory{#1}{#2}{#3}{relax}{0.5mu plus 0.25mu}{0.5mu plus 0.15mu}
\@bracketfactory{b#1}{#2}{#3}{big}{1mu plus 0.25mu minus 0.25mu}{0.6mu plus 0.15mu minus 0.15mu}
\@bracketfactory{bb#1}{#2}{#3}{Big}{2.4mu plus 0.8mu minus 0.8mu}{1.8mu plus 0.6mu minus 0.6mu}
\@bracketfactory{bbb#1}{#2}{#3}{bigg}{3.2mu plus 1mu minus 1mu}{2.4mu plus 0.75mu minus 0.75mu}
\@bracketfactory{bbbb#1}{#2}{#3}{Bigg}{4mu plus 1mu minus 1mu}{3mu plus 0.75mu minus 0.75mu}
}
\bracketfactory{clc}{\lbrace}{\rbrace}
\bracketfactory{clr}{(}{)}
\bracketfactory{cls}{[}{]}
\bracketfactory{abs}{\lvert}{\rvert}
\bracketfactory{norm}{\Vert}{\Vert}
\bracketfactory{floor}{\lfloor}{\rfloor}
\bracketfactory{ceil}{\lceil}{\rceil}
\bracketfactory{angle}{\langle}{\rangle}

%%%Define scr commands
\def\scrF{\mathscr{F}}
\def\scrG{\mathscr{G}}

%%% Define corresponding vertical bar commands

\def\calF{\mathcal{F}}
\def\calG{\mathcal{G}}
\def\ahalf{{\textstyle\frac12}}
\def\eq#1{\eqref{#1}}
\def\I{\mathrm{I}}
\def\Bi{\mathrm{Bi}}
\def\IP{\prob}
% \newcommand{\tunder}[1]{\underaccent{\tilde}{#1}}
\newcommand{\edgevar}[1]{[{#1}]_{\sim}}

% \newcommand{\edgevar}{[e]_{\sim}}
\def\ER{Erd\H{o}s-R\'enyi}
% \def\hKin{\hat{K}^{\mathrm{in}}}
% \def\hKout{\hat{K}^{\mathrm{out}}}
% \def\Kin{K^{\mathrm{in}}}
% \def\Kout{K^{\mathrm{out}}}

\makeatother


%Shortcuts for sets of numbers
\newcommand{\law}{\mathscr{L}}
\newcommand{\HH}{\mathscr{H}}
\newcommand{\D}{\mathbb{D}}
\newcommand{\Pro}{\mathbb{P}} 
\newcommand{\prob}{\Pro}
\DeclareMathOperator{\E}{\mathbb{E}}
% \newcommand{\E}{\E}
\newcommand{\mean}{\E}
\newcommand{\R}{\mathbb{R}}
\newcommand{\N}{\mathbb{N}}
\newcommand{\non}{\nonumber}
\newcommand{\Z}{\mathbb{Z}}
\newcommand{\C}{\mathbb{C}}
\newcommand{\G}{\mathbb{G}}
\newcommand{\LL}{\textbf{L}}
\DeclareMathOperator{\Var}{\mathrm{Var}}
\DeclareMathOperator{\var}{\mathrm{Var}}
\DeclareMathOperator{\cov}{\mathrm{Cov}}
\DeclareMathOperator{\bigo}{\mathrm{O}}
\newcommand{\lito}{\mathop{{}\mathrm{o}}\mathopen{}}
\newcommand{\K}{\textbf{Ker}}
\newcommand{\Id}{\textbf{Id}}
\newcommand{\card}{\mathop{\mathrm{card}}}

%%%%%%%%%%%%%
%%%%%Xia's commands
%%%%%%%%%%%%%
\usepackage{subfig}

\newcommand{\bn}{{\bf N}}
\newcommand{\bone}{{\bf 1}}
\newcommand{\bc}{{\mathds{C}}}
\newcommand{\bp}{{\bf P}}
\newcommand{\bh}{{\textbf H}}
\newcommand{\bi}{{\bf i}}
\newcommand{\bk}{{\bf k}}
\newcommand{\br}{{\bf R}}
\newcommand{\bq}{{\bf Q}}
\newcommand{\bs}{{\bf S}}
\newcommand{\bz}{{\bf Z}}
\newcommand{\bG}{\mbox{\boldmath$\Gamma$}}
\newcommand{\ba}{{\bm\alpha}}
\newcommand{\bl}{{\bm\lambda}}
\newcommand{\bmu}{{\bm\mu}}
\newcommand{\bx}{{\bm x}}
\newcommand{\bX}{{\bm X}}
\newcommand{\bY}{{\bm Y}}
\newcommand{\bZ}{{\bm Z}}
\newcommand{\cF}{{\cal F}}
\newcommand{\ca}{{\cal A}}
\newcommand{\cb}{{\cal B}}
\newcommand{\scrB}{{\mathscr B}}
\newcommand{\cc}{{\cal C}}
\newcommand{\cd}{{\cal D}}
\newcommand{\cf}{{\cal F}}
\newcommand{\cg}{{\cal G}}
\newcommand{\ch}{{\cal H}}
\newcommand{\ci}{{\cal I}}
\newcommand{\cl}{{\cal L}}
\newcommand{\cm}{{\cal M}}
\newcommand{\cn}{{\cal N}}
\newcommand{\cs}{{\cal S}}
\newcommand{\cu}{{\cal U}}
\newcommand{\cy}{{\cal Y}}
\newcommand{\dtvr}{{d_{\rm tv}}}%total variation for rvs
\newcommand{\dtv}{{d_{\rm TV}}}
\newcommand{\dk}{{d_{\rm K}}}
\newcommand{\dw}{{d_{\rm W}}}
\def\dkr{d_{\mathrm{KR}}}
\def\tg{{\tilde g}}
\def\tf{{\tilde f}}
\newcommand{\Pn}{{\rm Pn}}
\newcommand{\Po}{{\rm Po}}
\def\t#1{^{(#1)}}

\DeclareMathOperator*{\esssup}{ess\,sup}
\DeclareMathOperator*{\essinf}{ess\,inf}
% \newcommand{\Pn}{{\rm Pn}}

\def\Ref#1{(\ref{#1})}
\def\a{\alpha}
\def\b{\beta}
\def\s{\sigma}
\def\f{\phi}
\def\w{\omega}
%\def\l{\lambda}
\def\La{\Lambda}
\newcommand{\oF}{{\overline{F}}}
\def\p{\pi}
\def\tod{\buildrel{{\rm d}{\to}}}
\def\eqd{\stackrel{\rm d}{=}}
\def\siim{\sum_{i=1}^m}
\def\sjn{\sum_{j=1}^n}

% \def\tpj{\overline{\tau^+_j}}
% \def\tmj{\overline{\tau^-_j}}
% \def\tpi{\overline{\tau^+_i}}
% \def\tpi1{\overline{\tau^+_{i-1}}}
% \def\tmi{\overline{\tau^-_i}}
\def\stg{\stackrel{\mbox{\scriptsize st}}{\ge }}
\def\stl{\stackrel{\rm{\scriptsize st}}{\le }}

\newcommand*\rot[1]{\rotatebox{90}{#1}}

\makeatletter

% \newcommand{\qed}{\box}
%\newcommand{\qed}{\nopagebreak\hspace*{\fill}{\vrule width6pt height6ptdepth0pt}\par}


\newcommand{\blue}[1]{\textcolor{blue}{#1}}
\newcommand{\red}[1]{\textcolor{red}{#1}}
\newcommand{\yellow}[1]{\textcolor{yellow}{#1}}
\newcommand{\green}[1]{\textcolor{green}{#1}}
\newcommand{\cyan}[1]{\textcolor{cyan}{#1}}
\newcommand{\magenta}[1]{\textcolor{magenta}{#1}}

\def\ignore#1{}

\def\Comment#1{
    \marginpar{$\bullet$\quad{\tiny #1}}}

%%%%%definition of counters

\makeatletter
\newcommand*\labelcounter[2]{\begingroup
  \protected@edef\@currentlabel{\csname p@#1\endcsname\csname the#1\endcsname}%
  \label{#2}\endgroup}
\newcommand*\refsetcounter[2]{\setcounter{#1}{#2}%
  \protected@edef\@currentlabel{\csname p@#1\endcsname\csname the#1\endcsname}%
  }
\makeatother

\newcounter{rtaskno}
\DeclareRobustCommand{\rtask}[1]{%
   \refstepcounter{rtaskno}%
   \thertaskno\label{#1}}

\newcounter{con}%for constants
\newcommand{\const}{{\stepcounter{con}\rm\thecon}}
%to use the counter number, type \const and then immediately use \qcon{} so that the counter will be added by 1.
%Use \labelcounter{con}{eq:123} to extract the numeric number of the counter at this point

%%%%%end of definition of counters

%%%%%%%%%%%%%%
%end of Xia's commands
%%%%%%%%%%%%%%

\numberwithin{equation}{section}

\begin{document}
\title{\sc\bf\large{On Ergodicity of Spatial Birth--Death processes}}

\author{
Qingwei Liu\footnote{School of Mathematics and Statistics,
	The University of Melbourne,
	VIC 3010, Australia. E-mail: qingwei.liu@unimelb.edu.au.}\\ 
Aihua Xia\footnote{School of Mathematics and Statistics,
	The University of Melbourne,
	VIC 3010, Australia. E-mail: aihuaxia@unimelb.edu.au. Work supported in part by Australian Research Council Grant No DP220102666.}\\ {\it University of Melbourne}
	}

\def\parsedate #1:20#2#3#4#5#6#7#8\empty{20#2#3-#4#5-#6#7}
\def\moddate{\expandafter\parsedate\pdffilemoddate{\jobname.tex}\empty}
\date{\moddate}

\maketitle

\begin{abstract}

\end{abstract}


\vskip12pt \noindent\textit {Key words and phrases\/}: 

\vskip12pt \noindent\textit{AMS 2020 Subject Classification\/}: 


\section{Introduction}


\section{Preliminaries}
\subsection*{Point processes}
The carrier space $S$ of the point processes is complete and separable, equipped with metric $\varrho$ and corresponding Borel $\sigma$-field $\mathscr{B}(S)$.
Since a subset of $S$ can equivalently be represented as a non-negative integer-valued measure, we consider the space $\mathcal{N}$ of locally finite non-negative integer-valued measures on $(S,\mathscr{B}(S))$. 
%Let $\mathcal{N}_s$ denote the subspace of $\mathcal{N}$ of locally finite measures $\xi\in\mathcal{N}$ satisfying $\xi(\{x\})=0$ or $1$ for all $x\in S$.
Define $\mathscr{H}$ as the smallest $\sigma$-field on $\mathcal{N}$ such that the mappings $\mu\mapsto\mu(B)$ are measurable for all relatively compact sets $B$ in $\mathscr{B}(S)$~\cite[\S1.1]{kallenberg83}. Let $\mathcal{N}_{\rm f}\subseteq \mathcal{N}$ be the subspace of finite configurations on $(S,\mathscr{B}(S))$, and $\mathscr{H}_{\rm f}$ be the corresponding $\sigma$-algebra generated by the projection mappings.
%, and $\mathscr{H}_s{:=}\mathscr{H}\cap \mathcal{N}_s$. 
A point process $\mathcal{P}$ is a measurable mapping from an underlying probability space to $(\mathcal{N},\mathscr{H})$.
%A point process $\mathcal{P}$ is called {\it simple} if
%$$\prob\left(\mathcal{P}\in \mathcal{N}_s\right)=1.$$
Denote $\mathscr{P}(\mathcal{N})$ the space of all probability distributions on $(\mathcal{N},\mathscr{H})$.

For two configurations $\xi_1,\xi_2\in\mathcal{N}$, we denote $|\xi_1-\xi_2|$ as the total variation of the signed measure $\xi_1-\xi_2$, that is,
$|\xi_1-\xi_2|(B)=\xi^+(B)+\xi^-(B)$, where $\xi^+$ and $\xi^-$ represent the upper and lower variations of $\xi_1-\xi_2$, respectively, as described in the Jordan decomposition~\cite[p.~421]{B95}. Additionally, we define the norm $\|\xi_1-\xi_2\|:=\xi^+(S)+\xi^-(S)$.

Let $\mathcal{L}(\Xi)$ denote the distribution of a random element $\Xi$. %Let $\Pn(\theta)$ denote the Poisson distribution with parameter $\theta > 0$, and $\Po(\Lambda)$ as a Poisson point process with mean measure $\Lambda$.

We further denote
$$
\mathcal{N}_n:=\{\xi\in\mathcal{N}:\|\xi\|=n\},~~~
\mathcal{N}_{>n}:=\cup_{k=n+1}^{\infty}\mathcal{N}_k,
$$
and $\mathscr{H}_n:=\mathscr{H}\cap\mathcal{N}_n$, for $n\ge0$. 
Similarly, we write the $\N:=\{0,1,2,\dots\}$ and 
$$\N_{>n}:=\{n+1,n+2,\dots\}.$$
\section{Spatial birth--death processes}
\subsection{Construction}
A spatial birth--death process $\{X_t\}_{t\ge0}$, first introduced by \cite{Preston75}, is a Markov process defined on the configuration space $(\mathcal{N},\mathscr{H})$.
In the seminal paper \cite{Preston75}, the spatial birth--death process was constructed to allow only single births and single deaths. 
In this note, we extend the framework to consider a spatial birth--death process that permits multiple births but only single death.
Suppose the process evolves over time through the birth and death of individuals, with the mechanism governing this evolution specified by the following construction.
\begin{enumerate}
	\item At any time $t$, there are only finite individuals alive;
	\item Each time, if there is a birth, it would bring a cluster of individuals.

\item Suppose that there are $n$ individuals alive at time $t$, with locations 
$\xi = \{x_1, \dots, x_n\}$. For each $\ell \ge 1$, assume that there exist 
a measurable function $\lambda_\ell : \mathcal{N} \to [0, \infty)$ and 
a finite, symmetric measure $b_\ell(\xi, \cdot)$ on 
$(S^\ell, \mathscr{B}(S)^{\otimes \ell})$ such that:
\begin{enumerate}
  \item[i)] the probability that $\ell$ new individuals are born is 
  $\ell\, \lambda_\ell(\xi)$;
  \item[ii)] conditional on the event that $\ell$ new individuals are born, 
  for each $F \in \mathscr{B}(S)^{\otimes \ell}$ the probability that their 
  locations lie in $F$ during the interval $[t, t+s]$ is given by
  \[
    b_\ell(\xi, F)\, s + o(s).
  \]
\end{enumerate}
\item If there are $n$ individuals alive at time $t$, with locations $\xi=\{x_1,\dots,x_n\}$, then there exits a $\mathscr{B}(S)$-measurable function $\red{D(\xi\setminus x,\cdot)}: S\to[0,\infty)$ such that the probability of the individual at the point $x_i\in\xi$ dying between times $t$ and $t+s$ is $\red{D(\xi\setminus x_i,x_i)}s+o(s)$;
	\item The probability of more than one change in the process between times $t$ and $t+s$ is $o(s)$.
\end{enumerate}
More formally, we can construct the Markov process $\{X_t\}_{t\ge0}$ from a jump process $\{Y_n\}_{n\ge0}$. 
Let $\nu\in\mathscr{P}(\mathcal{N})$ be an initial distribution.
Let $D:\mathcal{N}\times S \to [0,\infty)$ be a bounded, measurable function. For each $\ell \ge 1$, let 
$\lambda_\ell : \mathcal{N} \to [0,\infty)$ be measurable and bounded, and let 
$b_\ell : \mathcal{N} \times \mathscr{B}(S)^{\otimes \ell} \to [0,\infty)$ be a kernel such that, for every 
$\xi \in \mathcal{N}$, the measure $b_\ell(\xi, \cdot)$ is symmetric on 
$(S^\ell, \mathscr{B}(S)^{\otimes \ell})$. Define
\begin{align}
    b(\xi) &:= \sum_{\ell=1}^{\infty} \ell\, \lambda_\ell(\xi)\, b_\ell(\xi, S^\ell),
\end{align}
for any $\xi \in \mathcal{N}$. Furthermore, define
\begin{align}
	d(\xi) &:= \sum_{x \in \xi} \red{D(\xi\setminus x, x)},
\end{align}
for all $\xi \neq \varnothing$, and set $d(\varnothing) := 0$.  
Finally, denote
\[
    \lambda(\xi) := b(\xi) + d(\xi).
\]

\begin{enumerate}
	\item $\{Y_0,Y_1,Y_2,\dots\}$ is a Markov chain in $\mathcal{N}$ satisfies that
\begin{enumerate}
	\item For any $M\in\mathscr{H}$, $\IP(Y_0\in M)=\nu(M)$;
	\item Given $Y_n=\xi$, then 
\begin{equation}
	Y_{n+1}:=\begin{cases}
		\xi\cup\{x_1,\dots,x_\ell\},\ \ \ &\mbox{w.p.}~\frac{1}{\lambda(\xi)}\ell\lambda_{\ell}(\xi)b_{\ell}(\xi,\mathrm{d}(x_1,\dots,x_{\ell}));\\
		\xi\setminus x,\ \ \ &\mbox{w.p.}~\frac{\red{D(\xi\setminus x,x)}}{\lambda(\xi)};
	\end{cases}
\end{equation}
\item That is $\{Y_n\}_{n\ge0}$ has initial distribution $\nu$ and the transition function
\begin{align}
	K(\xi,M):&=\IP(Y_{n+1}\in M|Y_n=\xi)\nonumber\\
		  &=\frac{1}{\lambda(\xi)}\sum_{\ell=1}^{\infty}\ell\lambda_{\ell}(\xi)\int_{S^{\ell}}\bone(\xi\cup\{x_1,\dots,x_\ell\}\in M)b_{\ell}(\xi,\mathrm{d}(x_1,\dots,x_{\ell}))\nonumber\\
		  &\ \ \ \ +\frac{1}{\lambda(\xi)}\sum_{x\in\xi}D(\xi\setminus x,x)\bone(\xi\setminus x\in M).\label{trans1}
\end{align}
\end{enumerate}
\item Let $\{T_0,T_1,T_2,\dots\}$ be a sequence of i.i.d. exponential random variables with parameter $1$ and indepdent of $\{Y_n\}_{n\ge0}$.
\item Define 
\begin{equation}
	X_t:=\begin{cases}Y_0,\ \ \ &0\le t \le \frac{T_0}{\lambda(Y_0)};\\
		Y_k,\ \ \ &\sum_{j=0}^{k-1}\frac{T_j}{\lambda(Y_j)}\le t<\sum_{j=0}^k\frac{T_j}{\lambda(Y_j)}.
	\end{cases}
\end{equation}
\end{enumerate}
%Note that $\lambda(\cdot)$ is the intensity of the Markov jump process.
Denote $\lambda:=\sup_{\xi\in\mathcal{N}}\lambda(\xi)$. To avoid trivialities, assume that $\lambda>0$. 


Given $X_t=\xi$ with $\xi\in\mathcal{N}$ and $t\ge0$, then the waiting time to the next jump has an exponential distribution with expectation $1/\lambda(\xi)$ and is independent of the past history. The subsequent jump leads either to a death of an individual located at point $x\in\xi$ with probability $\frac{\red{D(\xi\setminus x,x)}}{\lambda(\xi)}$, or a birth of $\ell$ individuals, distributed according to $b_{\ell}(\xi,\cdot)$.  

The generator $\mathscr{A}$ of $\{X_t\}_{t\ge0}$ is given by
\begin{equation}
	(\mathscr{A}f)(\xi)=\lambda(\xi)\int(f(\zeta)-f(\xi))K(\xi,\mathrm{d}\zeta),
\end{equation}
which can also be rewrited as 

\begin{align}
	(\mathscr{A}f)(\xi)&=\sum_{\ell=1}^\infty \ell \lambda_\ell(\xi)\int_{S^\ell}\left(f(\xi\cup\{x_1,\dots,x_\ell\})-f(\xi)\right)b_\ell(\xi,\mathrm{d}(x_1\cdots x_\ell))\nonumber\\
			   &\ \	+\int_{S}\red{D(\xi\setminus x,x)}\left(f(\xi\setminus x)-f(\xi)\right)\xi(\mathrm{d}x).
\end{align}
for any bounded measurable function $f:\mathcal{N}\to[0,\infty)$.
We call the functions $b: \mathcal{N}\to[0,\infty)$ the birth rate coefficient and $d:\mathcal{N}\to[0,\infty)$ the death rate coefficient. 

Consequently, the semigroup $\{T(t)\}$ generated by $\mathscr{A}$ is given by 
\begin{equation}
	T(t):=\sum_{k=0}^{\infty}\frac{(t\mathscr{A})^k}{k!}=\exp(t\mathscr{A}).
\end{equation}

Note this spatical birth--death process allows multiple birth and single death. The case that only single birth and single death are allowed, is often referred as {\em Glauber dynamics in continuum}, see \cite{KL05} and \cite{moller03book}. 


\subsection{Uniqueness and backwards equations}

To define such a process analytically, we also need to give its transition semigroup, i.e., for $t\ge0$, we need to define $Q_t:\mathcal{N}\times\mathscr{H}\to[0,1]$ such that 
$$Q_t(\xi,M):=\IP(X_t\in M|X_0=\xi),$$
for any $\xi\in\mathcal{N}$ and $M\in\mathscr{H}$. 
Let $Q_t$ be the solution of Kolmogorov's backwards equations
\begin{equation}\label{backweq-1}
	\frac{\partial}{\partial t}Q_t(\xi,M)=-\lambda(\xi)Q_t(\xi,M)+\lambda(\xi)\int Q_t(\zeta,M)K(\xi,\mathrm{d}\zeta),
\end{equation}
with the initial conditions 
\begin{equation}\label{backweq-2}
Q_0(\xi,M)=\delta_{\xi}(M).
\end{equation}
Let $Q_t^{(n)}(\xi,M)$ be the probability of a transition from state $\xi$ to $M$ in time $t$ with at most $n$ jumps. 
Let 
\begin{equation}
	Q_t^{(\infty)}(\xi,M):=\lim_{n\to\infty}Q_t^{(n)}(\xi,M).
\end{equation}
$Q_t^{(\infty)}$ is the minimal solution of the backwards equations \eqref{backweq-1} and \eqref{backweq-2}.
Recall the fact about the minimal solution: \eqref{backweq-1} and \eqref{backweq-2} have a unique solution if and only if $Q_t^{(\infty)}(\xi,\mathcal{N})=1$ for all $t\ge0$ and $\xi\in\mathcal{N}$. 
For $n\ge0$, let $\leftsup{(n)}{Q_t}(\xi,M)$ be the probability of a transition from $\xi$ to $M$ in time $t$ without having entered $\mathcal{N}_{>n}:=\cup_{k=n+1}^{\infty}\mathcal{N}_k$, and 
$$
\leftsup{(\infty)}{Q_t}(\xi,M):=\lim_{n\to\infty}\leftsup{(n)}{Q_t}(\xi,M).
$$
\begin{lma}\label{lma1}
	For any $\xi\in\mathcal{N}$, $M\in\mathscr{H}$, it holds 
\begin{equation}
	\leftsup{(\infty)}{Q_t}(\xi,M)=Q_t^{(\infty)}(\xi,M).
\end{equation}
\end{lma}
\begin{proof}
	\textcolor{cyan}{It is easy to check that $\leftsup{(\infty)}{Q_t}$ satisfies \eqref{backweq-1} and \eqref{backweq-2}}, and that for any $n$, $\leftsup{(n)}{Q_t}\le Q_t^{(\infty)}$. 
\end{proof}


\subsection{Construction of a coupled single death process}
Following the idea in \cite{Preston75}, it is often convenient to transform the analysis of a spatial birth--death process into that of a coupled simple birth--death process.

Set for any $n\ge1$, 
\begin{align}
	q_{n,n-1}&:=\inf_{\xi:\|\xi\|=n}d(\xi)=\inf_{\xi:\|\xi\|=n}\sum_{x\in\xi}\red{D(\xi\setminus x,x)};\label{qmatrix1}\\
	q_{i,j}&:=0,~~~\forall~j<i-1;\\
	q_{n,n+\ell}&:=\sup_{\xi:\|\xi\|=n}\ell\lambda_\ell(\xi) b_{\ell}(\xi,S^{\ell}),~~~\forall~\ell\ge1;\\
	q_{n,n}&:=-q_{n,n-1}-\sum_{\ell=1}^{\infty}q_{n,n+\ell}.\label{qmatrix4}
\end{align}
Notationally, we also set $q_n:=-q_{n,n}$. 

Let $\{Z_t\}_{t\ge0}$ denote a single death process on $\N:=\{0,1,2,\dots\}$ with $Q$-matrix $Q=(q_{i,j})$.
The single-death process $\{Z_t\}_{t \ge 0}$ is a simple birth--death process that allows for multiple births but only single deaths. Such a process is also referred to as a \emph{downwardly skip-free process} \cite{CPZC05}.
Suppose that a single death process is irreducible and recurrent, \cite[Theorem~1.1]{Zhang18} establishes the necessary and sufficient conditions for the process to be ergodic and strongly ergodic, respectively.

In order to prove our main theorm, we need to compare our spatial birth-death process with the single death process with $Q$-matrix given in \eqref{qmatrix1}-\eqref{qmatrix4}.
\begin{thm}\label{thm-uni1}
	Suppose that the backwards equations for the single death process $\{Z_t\}_{t\ge0}$ have a unique solution $q_t$. Then the backwards equations corresponding to $\lambda$ and $K$ have a unique solution $Q_t$.
\end{thm}
\begin{proof}
Let us first formally restate the definitions of $\leftsup{(b)}{Q}_t(\xi,M)$ and $\leftsup{(m)}{q}_t(n,K)$.
\begin{enumerate}
	\item Let $\leftsup{(m)}{Q}_t(\xi,M)$ denote the probability of the spatial birth-death process $\{X_t\}_{t\ge0}$ transitioning from state $\xi$ to a set of states $M$ in time $t$, given the process has not entered the set $\mathcal{N}_{>m}$ during the interval $[0, t]$, for any $\xi\in\mathcal{N}$ and $M\in\mathscr{H}$.
	\item Let $\leftsup{(k)}{q}_t(n,K)$ denote the probability of the single death process $\{Z_t\}_{t\ge0}$ transitioning from state $n$ to $K$ in time $t$, given the process has not entered the set $\N_{>k}$ during the interval $[0, t]$, for any $n\in\N$ and $K\subseteq\N$.
\end{enumerate}
We then construct a new jump process $\{(\tilde{X}_t,\tilde{Z}_t)\}_{t\ge0}$ with state space $(\mathcal{N}\times\N,\mathscr{H}\times 2^{\N})$.
The coupled process $\{(\tilde{X}_t,\tilde{Z}_t)\}_{t\ge0}$ is defined by its intensity $\tilde{\lambda}$ and transition function $\tilde{K}$ as follows:
\begin{equation}
	\tilde{\lambda}(\xi,n):=
\begin{cases}
	b(\xi)+d(\xi)+q_n,\ \ \ &\mbox{if}~\xi\in\mathcal{N}_m,~m\ne n,\\
	d(\xi)+\sum_{\ell=1}^{\infty}q_{n,n+\ell},\ \ \ &\mbox{if}~\xi\in\mathcal{N}_n,
\end{cases}
\end{equation}
and for any $\xi\in\mathcal{N}$, $M\in\mathscr{H}$, 
\begin{enumerate}
	\item when $\xi\in\mathcal{N}_m$, and $m\ne n$, 
\begin{align}
	\tilde{K}(\xi,n;M\times\{n\})&=\frac{\lambda(\xi)}{\tilde{\lambda}(\xi,n)}K(\xi,M),\nonumber\\
	\tilde{K}(\xi,n;\{\xi\}\times\{n+\ell\})&=\frac{q_{n,n+\ell}}{\tilde{\lambda}(\xi,n)},~~~\ell\ge1,\\
	\tilde{K}(\xi,n;\{\xi\}\times\{n-1\})&=\frac{q_{n,n-1}}{\tilde{\lambda}(\xi,n)};\nonumber
\end{align}
\item when $\xi\in\mathcal{N}_n$,
\begin{align}
	\tilde{K}(\xi,n;M_{n+\ell}\times\{n+\ell\})&=\frac{\ell b_{\ell}(\xi,S^{\ell})}{\tilde{\lambda}(\xi,n)}K(\xi,M_{n+\ell}),\ \ \ \ell\ge1,\nonumber\\
	\tilde{K}(\xi,n;\{\xi\}\times\{n+\ell\})&=\frac{q_{n,n+\ell}-\ell b_{\ell}(\xi,S^{\ell})}{\tilde{\lambda}(\xi,n)},\ \ \ \ell\ge1,\nonumber\\
	\tilde{K}(\xi,n;M_{n-1}\times\{n-1\})&=\frac{q_{n,n-1}}{\tilde{\lambda}(\xi,n)}K(\xi,M_{n-1}),\\
	\tilde{K}(\xi,n;M_{n-1}\times\{n\})&=\frac{d(\xi)-q_{n,n-1}}{\tilde{\lambda}(\xi,n)}K(\xi,M_{n-1}),\nonumber
\end{align}
where $M_k:=M\cap\mathcal{N}_k$ for all $k\ge0$.
\end{enumerate}
Similarly, let $\leftsup{(m,k)}{\tilde{Q}}(\xi,n;F)$ denote the probability of a transition from $(\xi,n)$ to $F\subseteq\mathcal{N}\times\N$ in time $t$ without having entered 
$$
\{(\zeta,j)\in\mathcal{N}\times\N:~\zeta\in\mathcal{N}_{>m} ~\mbox{or}~ j>k\}
$$
From the construction of the coupled process $(\tilde{X}_t,\tilde{Z}_t)$, it is easy to check that it satisfies the following properties:
\begin{enumerate}
	\item[i)] For any $\xi\in\mathcal{N}$, $K\subseteq \N$,
\begin{equation}
	\tilde{Q}_t(\xi,n;\mathcal{N}\times K)=q_t(n,K);
\end{equation}
\item[ii)] For any $\xi\in\mathcal{N}$, $M\in\mathscr{H}$,  $\tilde{Q}_t(\xi,n;M\times\N)$ is independent of $n\in\N$ and if $Q_t(\xi,M)=\tilde{Q}_t(\xi,n;M\times\N)$ then $Q_t$ is the minimal solution of the backwards equations for the original process $\{X_t\}_{t\ge0}$;
\item[iii)] If $\xi\in\mathcal{N}_m$ and $m\le n$, then 
	$$
	\tilde{Q}_t(\xi,n;\Gamma)=0
	$$
	for all $t\ge0$, where
	$$\Gamma=\{(\zeta,\tilde{m})\in\mathcal{N}\times \N:\zeta\in\mathcal{N}_k,~~\mbox{and}~ k>\tilde{m}\}.$$
\end{enumerate}
Furthermore, for any $\xi\in\mathcal{N}_j$ with $j\le n$, we have 
\begin{equation}
	\leftsup{(m,k)}{\tilde{Q}}_t(\xi,n;F)=\leftsup{(\tilde{m},k)}{\tilde{Q}}_t(\xi,n;F)
\end{equation}
for any $m,\tilde{m}>k$,
and
\begin{equation}
	\leftsup{(m,k)}{\tilde{Q}}_t(\xi,n;\mathcal{N}\times K)=\leftsup{(k)}{q}_t(n,K)
\end{equation}
for any $m>k$.
Therefore, according to Lemma~\ref{lma1}, we have, for any $\xi\in\mathcal{N}_j$ with $j\le n$
\begin{equation}
	q^{(\infty)}_t(n,K)=\leftsup{(\infty)}{q}_t(n,K)=\lim_{k\to\infty}\leftsup{(k+1,k)}{\tilde{Q}}_t(\xi,n;\mathcal{N}\times K),
\end{equation}
and again by Lemma~\ref{lma1}, %we have, for any $\xi\in\mathcal{N}_j$ with $j\le n$,
\begin{equation}
  \lim_{k\to\infty}\leftsup{(k+1,k)}{\tilde{Q}}_t(\xi,n;\mathcal{N}\times K)
  =\tilde{Q}^{(\infty)}_t(\xi,n;\mathcal{N}\times K).
\end{equation}
This leads to that 
for any $\xi\in\mathcal{N}_j$ with $j\le n$,
\begin{equation}
	q^{(\infty)}_t(n,K)=\tilde{Q}^{(\infty)}_t(\xi,n;\mathcal{N}\times K),
\end{equation}
and in particular, 
\begin{equation}
	\tilde{Q}^{(\infty)}_t(\xi,n;\mathcal{N}\times\N)=q^{(\infty)}_t(n,\N)=1.
\end{equation}
In the other direction, 
for any $m\ge0$
\begin{equation}
	\tilde{Q}^{(m)}_t(\xi,n;\mathcal{N}\times\N)\le Q^{(m)}_t(\xi,\mathcal{N}),
\end{equation}
thus 
\begin{equation}
	Q^{(\infty)}_t(\xi,\mathcal{N})=1.
\end{equation}
Therefore, the backwards equations \eqref{backweq-1} and \eqref{backweq-2} corresponding to $\lambda$ and $K$ has a unique solution.
\end{proof}


\section{Main theorem}
\subsection{Foster-Lyapunov condition}\label{sec:foster}

According to \cite[Theorem~3]{Tweedie81}, if $Q$ is regular, then $\{Z_t\}_{t\ge0}$ is 
\begin{enumerate}
	\item ergodic if and only if there is an $N\ge0$ and a finite solution $y_i\ge0$ to the inequality
		\begin{equation}\label{lyapunov-ineq}
	\sum_{k}q_{i,k}y_k\le -\lambda y_i-1,~~~i>N,
\end{equation}
and $y_i=0$ for $i=0,1,\dots,N$, with $\lambda=0$, such that 
\begin{equation}\label{cond-2}
	\sum_{k}\frac{(1-\delta_{i,k})q_{i,k}}{q_i}y_k<\infty,~~~i\le N;
\end{equation}
\item exponentially ergodic if and only if there is an $N\ge0$ and a finite solution $y_i\ge0$ to the inequalities \eqref{lyapunov-ineq} for some $\lambda>0$ with $\lambda<q_i$ for all $i$ such that \eqref{cond-2} holds.
\end{enumerate}

We can construct a specific test sequence $y_i$ that grows linearly with $i$ to test for both standard and exponential ergodicity. For the simplest setting, we take $N=0$ and $y_i=i$ for $i\ge0$, and $y_0=0$. Therefore, for $i>0$, 
\begin{align*}
	\sum_{k}q_{i,k}y_k=\sum_k q_{i,k}k&=\sum_{k>i}q_{i,k}k+q_{i,i-1}(i-1)+q_{i,i}i\\
					  &=\sum_{k>i}q_{i,k}k+q_{i,i-1}(i-1)-\left(\sum_{k>i}q_{i,k}+q_{i,i-1}\right)i\\
					  &=\sum_{\ell=1}^\infty q_{i,i+\ell}\ell-q_{i,i-1}.
\end{align*}
Let
\begin{equation}
	\Delta(i) := \sum_{\ell=1}^{\infty} \ell q_{i,i+\ell} - q_{i,i-1},
\end{equation}
for all $i\ge1$. We have the following sufficient conditions from Tweedie's criteria.

\begin{thm}
Suppose $Q$ is regular, and it holds that   
\begin{equation}
	\sum_{\ell=1}\ell q_{0,\ell}<\infty.
\end{equation}
Then $\{Z_t\}_{t\ge0}$ is ergodic if 
\begin{equation}
  \Delta(i)\le -1,
\end{equation}
for all $i\ge1$. \footnote{\red{Essentially, we only need 
$
\sup_{i>0}\Delta(i)<0.
$
}}
Moreover, $\{Z_t\}_{t\ge0}$ is exponentially ergodic if 
\begin{equation}
  \Delta(i)\le -\lambda i-1,
\end{equation}
for all $i\ge1$ with some $0<\lambda<q_i$, for all $i\ge1$.
\end{thm}

\subsection{An alternative approach}
Denote 
\begin{equation}
	q_n^{(k)}:=\sum_{j=k}^{\infty}q_{n,j},~~~k>n\ge0,
\end{equation}
and
\begin{equation}
	G_i^{(i)}:=1,~~~G_n^{(i)}:=\frac1{q_{n,n-1}}\sum_{k=n+1}^iq_n^{(k)}G_k^{(i)},~~~1\le n<i.
\end{equation}
\begin{thm}
	Assume that the single death process $\{Z_t\}_{t\ge0}$ with $Q$-matrix given in \eqref{qmatrix1}-\eqref{qmatrix4} is irreducible and the recurrent.
	If 
	\begin{equation}\label{ergod-1}
	\sum_{k\ge0}q_0^{(k)}\sum_{\ell\ge k}\frac{G_k^{(\ell)}}{q_{\ell,\ell-1}}<\infty,
\end{equation}
then the spatial birth-death process $\{X_t\}_{t\ge0}$ is ergodic. 
\end{thm}
\begin{proof}
	According to \cite[Theorem~1.1]{Zhang18}, when \eqref{ergod-1} holds, the single death process $\{Z_t\}_{t\ge0}$ is ergodic. It follows from Theorem~\ref{thm-uni1} that $\varnothing$ is an ergodic state for $\{X_t\}_{t\ge0}$. Thus, by the renewal theorem \cite[p.~379]{Feller71}, we get $\lim_{t\to\infty}Q_t(\varnothing,M)$ exists, for any $M\in\mathscr{H}$. Furthermore, Theorem~\ref{thm-uni1} ensures that for any $\xi\in\mathcal{N}$, a passage from $\xi$ to state $\varnothing$ is certain and the expected passage time is finite. By the renewal theorem, we get $\lim_{t\to\infty}Q_t(\xi,M)$ exists for any $\xi\in\mathcal{N}$ and $M\in\mathscr{H}$.
\end{proof}

\subsection{Stationary distribution}
The Foster-Lyapunov condition in Section~\ref{sec:foster} ensures the spatial multiple birth process to be ergodic. 

In the current section, we further study a (potential) stationary distribution for this process. 

Throughout this note, we assume the carrier space $S$ is a compact set of Euclid space $\R^d$ and we take $\rho$ be the Lebesgue measure on $\R^d$. 
\begin{defi}[Reduced Campbell measure and Papangelou intensity]
Let $\mathcal{P}$ be a point process on $S$ with law $\mu := \mathscr{L}(\mathcal{P})$ on $(\mathcal{N},\mathscr{H})$.
The reduced Campbell measure of $\mathcal{P}$ is the measure $C_{\mathcal{P}}$ on $(S\times \mathcal{N},\, \mathscr{B}(S)\otimes \mathscr{H})$
defined by
\begin{equation}\label{eq:reduced-campbell}
    C_{\mathcal{P}}(A)
    :=
    \int_{\mathcal{N}}
    \sum_{x\in \xi}
    \mathbf{1}_A\bigl(x,\xi\setminus x\bigr)\,
    \mu(\mathrm{d}\xi),
    \qquad A\in \mathscr{B}(S)\otimes \mathscr{H}.
\end{equation}
The reduced compound Campbell measure of $\mathcal{P}$ is the measure $\widehat{C}_{\mathcal{P}}$ on $(\mathcal{N}_{\rm f}\times \mathcal{N},\, \mathscr{H}_{\rm f}\otimes \mathscr{H})$
defined by
\begin{equation}\label{eq:reduced-compound-campbell}
    \widehat{C}_{\mathcal{P}}(B):=\int_{\mathcal{N}}
    \sum_{\alpha\in\mathcal{N}_{\rm f},\,\alpha\subset \xi}
    \mathbf{1}_B\bigl(\alpha,\xi\setminus \alpha\bigr)\,
    \mu(\mathrm{d}\xi),
    \qquad B\in \mathscr{H}_{\rm f}\otimes \mathscr{H}.
\end{equation}
\end{defi}

\begin{defi}
We say that $\mathcal{P}$ satisfies condition $(\Sigma_\rho)$ if
\[
    C_{\mathcal{P}} \ll \rho\otimes \mu,
\]
where $\rho$ denotes Lebesgue measure on $\R^d$.
In this case, any Radon--Nikodym derivative
\begin{equation}\label{eq:papangelou-def}
	c(x,\xi):=\frac{\mathrm{d} C_{\mathcal{P}}}{\mathrm{d}(\rho\otimes \mu)}(x,\xi),
    \qquad (x,\xi)\in S\times \mathcal{N},
\end{equation}
of $C_{\mathcal{P}}$ with respect to $\rho\otimes\mu$
is called a \emph{Papangelou (conditional) intensity} (or \emph{Papangelou density}) of $\mathcal{P}$.
\end{defi}
More explicitly, $c$ is a Papangelou intensity of $\mathcal{P}$ if
\begin{equation}
	\int\sum_{x\in\xi}f(x,\xi\setminus x)\mu(\mathrm{d}\xi)=\int_S\int_{\mathcal{N}}c(x,\xi)f(x,\xi)\mu(\mathrm{d}\xi)\rho(\mathrm{d}x)
\end{equation}
for all measurable function $f:S\times \mathcal{N}\to\R_+$, see more details in \cite{GY05}. 

Let $L$ denote the Lebesgue--Poisson measure on $(\mathcal{N}_{\rm f},\mathscr{H}_{\rm f})$, defined by the identity
\begin{equation}\label{eq:LP-measure}
\int_{\mathcal{N}_{\rm f}}f(\alpha)L(\mathrm{d}\alpha)=
\sum_{n=0}^{\infty}\frac{1}{n!}\int_{S^n}f\bigl(\{x_1,\ldots,x_n\}\bigr)\rho(\mathrm{d}x_1)\cdots\rho(\mathrm{d}x_n)
\end{equation}
for any measurable function $f:\mathcal{N}_{\rm f}\to\R_+$.
According to \cite[Remark~2.5]{GY05}, for any point process $\mathcal{P}$ satisfying condition $(\Sigma_\rho)$, the reduced compound Campbell measure $\widehat{C}_{\mathcal{P}}$ is absolutely continuous with respect to $L\otimes \mu$ with a Radon--Nikodym density $\widehat{c}$ satisfying $\widehat{c}(\varnothing,\xi)=1$ and 
\begin{equation}\label{eq:compound-papangelou}
	\widehat{c}(\alpha,\xi)=c(x_1,\xi)\prod_{i=2}^nc(x_i,\{x_1,\ldots,x_{i-1}\}\cup\xi),
\end{equation}
where $\alpha=\{x_1,\ldots,x_n\}$.
%%%%%%%%%%%%%%%%%%

\begin{prop}\label{prop:papangelou-density-ratio}
Let $\Pi_\rho$ denote the law of the Poisson point process on $S$ with intensity measure $\rho$. Suppose that a probability measure $\pi$ on $(\mathcal{N},\mathscr{H})$ is absolutely continuous with respect to $\Pi_\rho$, and write
$$
    \pi(\mathrm{d}\xi)=p(\xi)\,\Pi_\rho(\mathrm{d}\xi)
$$
for some measurable function $p:\mathcal{N}\to[0,\infty)$.
Then $\pi$ satisfies condition $(\Sigma_\rho)$, and its Papangelou intensity is given by
\begin{equation}\label{eq:papangelou-density-ratio}
    c(x,\xi)=\frac{p(\xi+\delta_x)}{p(\xi)},
    \qquad \rho\otimes\pi\text{-a.e. }(x,\xi).
\end{equation}
Moreover, the corresponding compound Papangelou intensity satisfies
\begin{equation}\label{eq:compound-papangelou-density-ratio}
    \widehat{c}(\alpha,\xi)=\frac{p(\xi+\alpha)}{p(\xi)},
    \qquad L\otimes\pi\text{-a.e. }(\alpha,\xi),
\end{equation}
where we adopt the convention $0/0:=0$.
\end{prop}
\begin{proof}
	For any measurable function $f:S\times\mathcal{N}\to\R_+$, by the Mecke identity for the Poisson point process law $\Pi_\rho$
\begin{align}
	\int\sum_{x\in\xi}f(x,\xi\setminus x)\pi(\mathrm{d}\xi)&=\int\sum_{x\in\xi}f(x,\xi\setminus x)p(\xi)\Pi_\rho(\mathrm{d}\xi)\nonumber\\
									&=\int_S\int f(x,\xi)p(\xi+\delta_x)\Pi_\rho(\mathrm{d}\xi)\rho(\mathrm{d}x)\nonumber\\
									&=\int_S\int f(x,\xi)\frac{p(\xi+\delta_x)}{p(\xi)}p(\xi)\Pi_\rho(\mathrm{d}\xi)\rho(\mathrm{d}x)\nonumber\\
									&=\int_S\int f(x,\xi)\frac{p(\xi+\delta_x)}{p(\xi)}\pi(\mathrm{d}\xi)\rho(\mathrm{d}x),
\end{align}
with the convention $0/0:=0$.
Hence $\pi$ satisfies condition $(\Sigma_\rho)$, and
$$
    c(x,\xi)=\frac{p(\xi+\delta_x)}{p(\xi)},
    \qquad \rho\otimes\pi\text{-a.e. }(x,\xi).
$$
From \eqref{eq:compound-papangelou} and \eqref{eq:papangelou-density-ratio}, the identity \eqref{eq:compound-papangelou-density-ratio} follows from a telescopic argument.
\end{proof}


Assume that the spatial birth-death process $\{X_t\}_{t\ge0}$ admits an invariant probability measure $\pi$.

 %%%%%%%%%%%%%%%%%%%%%%%%%%%%%
%%%%%%% References   %%%%%%%%
%%%%%%%%%%%%%%%%%%%%%%%%%%%%%

\def\ac{{Academic Press}~}
\def\aap{{Adv. Appl. Prob.}~}
\def\ap{{Ann. Probab.}~}
\def\anap{{Ann. Appl. Probab.}~}
\def\eljp{{\it Electron.\ J.~Probab.\/}~} 
\def\jap{{J. Appl. Probab.}~}
\def\jws{{John Wiley~$\&$ Sons}~}
\def\ny{{New York}~}
\def\ptrf{{Probab. Theory Related Fields}~}
\def\sp{{Springer}~}
\def\spa{{Stochastic Process. Appl.}~}
\def\sv{{Springer-Verlag}~}
\def\tpa{{Theory Probab. Appl.}~}
\def\zw{{Z. Wahrsch. Verw. Gebiete}~}

% \bibliographystyle{apalike}
% \bibliography{ref}

\begin{thebibliography}{}


\bibitem[Barbour and Brown~(1992)]{BB92}
\newblock Barbour, A.~D. and Brown,~T.~C.~(1992). 
\newblock Stein's method and point process approximation. \newblock {\em \spa}~\textbf{43}, 9--31.

\bibitem[Billingsley~(1995)]{B95}
\newblock Billingsley, P.~(1995).
\newblock {\em Probability and Measure}.  
\newblock Third Edition. John Wiley \& Sons. 

\bibitem[Chen {\it et al}~(2005)]{CPZC05} 
\newblock Chen, A., Pollett, P., Zhang, H., and Cairns, B.~(2005).
\newblock Uniqueness criteria for continuous-time Markov chains with general transition structures.
\newblock {\em \aap}~\textbf{37}, 1056--1074.

\bibitem[Chen~(1975)]{Chen75} Chen, L. H. Y.~(1975).
\newblock Poisson approximation for dependent trials. 
\newblock {\em \ap}\textbf{3}, 534--545. 

\bibitem[Chen and Xia~(2004)]{CX04} Chen, L. H. Y  and Xia, A.~(2004) 
\newblock Stein's method, Palm theory and Poisson process approximation. 
\newblock {\em \ap}\textbf{32}, 2545--2569.

\bibitem[Chiu {\it et al}~(2013)]{chiu13}
\newblock Chiu, S. N., Stoyan, D., Kendall, W. S. and Mecke, J.~(2013)
\newblock {\em Stochastic geometry and its applications}. 
\newblock John Wiley \& Sons.

\bibitem[Daley and Vere-Jones~(2003)]{DV03}
Daley, D. J. and Vere-Jones, D.~(2003). 
\newblock {\em An introduction to the theory of point processes, Vol. I: Probability and its Applications}.  
\newblock Second edition, Springer-Verlag.

\bibitem[Daley and Vere-Jones~(2008)]{DV08}
Daley, D. J. and Vere-Jones, D.~(2008). 
\newblock {\em An introduction to the theory of point processes, Vol. II: General Theory and Structure}.  
\newblock Second edition, Springer-Verlag.

\bibitem[Ethier and Kurtz~(2009)]{EK09}
Ethier, S. N. and Kurtz T. G.~(2009).
\newblock {\em Markov processes: Characterization and Convergence}.
\newblock John Wiley \& Sons.

\bibitem[Feller~(1971)]{Feller71}
Feller, W.~(1971). 
\newblock {\em An Introduction to Probability Theory and Its Applications.\/}, Vol.~2, 2nd ed., 
\newblock John Wiley \& Sons, Inc.

\bibitem[Flint {\it et al.}~(2022)]{FX22}Flint, I., Golding, N., Vesk, P., Wang, Y. and Xia, A.~(2022). 
\newblock The saturated pairwise interaction Gibbs point process as a joint species distribution model. 
\newblock {\em J. R. Stat. Soc. Ser. C. Appl. Stat.}~\textbf{71}, 1721--1752.

\bibitem[Georgii and Yoo~(2005)]{GY05}Georgii, HO., and Yoo, H.J.~(2005).
\newblock Conditional Intensity and Gibbsianness of Determinantal Point Processes.
\newblock {\em J. Stat. Phys.}~\textbf{118}, 55--84.
\bibitem[Kallenberg~(2021)]{kallenberg21}
Kallenberg, O.~(2021).
\newblock {\em Foundations of Modern Probability}.
\newblock Third Edition. Springer Cham.

\bibitem[Kallenberg~(1983)]{kallenberg83}
Kallenberg, O.~(1983).
\newblock {\em Random Measures}.
\newblock De Gruyter, Berlin, Boston.

\bibitem[Kondratiev and Lytvynov~(2005)]{KL05}
Kondratiev, Y. and Lytvynov, E.~(2005)
\newblock Glauber dynamics of continuous particle systems
\newblock {\em Ann. Inst. Henri Poincar\'e Probab. Stat.}~\textbf{41}, 685--702. 

\bibitem[Moller and Waagepetersen~(2003)]{moller03book}
Moller, J. and Waagepetersen, R.P.~(2003)
\newblock {\rm Statistical inference and simulation for spatial point processes}.
\newblock Chapman and Hall/CRC.

\bibitem[Preston~(1975)]{Preston75}
Preston, C.~(1975).
\newblock Spatial birth and death processes.
\newblock {\em \aap}~\textbf{7}, 371--391.

\bibitem[Tweedie~(1981)]{Tweedie81}
Tweedie, R. L.~(1981).
\newblock Criteria for Ergodicity, Exponential Ergodicity and Strong Ergodicity of Markov Processes.
\newblock {\em \jap}~\textbf{18}, 122--130.

\bibitem[Zhang~(2018)]{Zhang18}
Zhang, Y.~(2018)
\newblock Criteria on ergodicity and strong ergodicity of single death processes.
\newblock {\em Front. Math. China}~\textbf{13}, 1215--1243.

  \end{thebibliography}
\end{document}









